% 基本信息--------------------------------------------------------------------------------
% 指定论文类型0:毕业论文  1/2:专业实践1/2
\newcommand{\reportStyle}{0}

% 在这里填写你的论文中文题目
\newcommand{\thesisTitle}{基于区块链的去中心化音乐社区平台设计与实现}
% 在这里填写你的论文英文题目
\newcommand{\thesisTitleEN}{Design and Implementation of a Blockchain-Based Decentralized Music Community Platform}


% 若有副标题，则运行下一行的代码，若无副标题，则将下一行注释掉（在\haveSub{}最前面添加 % 号）
% \haveSub{}

% 在这里填写你的论文中文副标题（没有只需注释掉\haveSub{}即可）
\newcommand{\thesisSubTitle}{融合链上音乐资源检索与排序模块}
% 在这里填写你的论文英文副标题（没有只需注释掉\haveSub{}即可）
\newcommand{\thesisSubTitleEN}{With an On-chain Music Resource Retrieval and Ranking Module}

% 在这里填写你的相关信息
\newcommand{\deptName}{信息技术与人工智能学院}
\newcommand{\majorName}{软件工程}
\newcommand{\yourName}{陈普阳}
\newcommand{\yourStudentID}{220110790202}
\newcommand{\mentorName}{陈伟锋}
\newcommand{\className}{22软件2班}
\newcommand{\Today}{2026年5月}
% 基本信息--------------------------------------------------------------------------------

% 中英文摘要与关键词
\newcommand{\abstractCN}{
    随着数字音乐发行方式和区块链技术的发展，传统音乐平台在创作者发行流程、版权凭证、收益分配和资源访问控制等方面仍存在一定的中心化依赖。针对上述问题，本文设计并实现了一个面向音乐社区场景的链上音乐发行与播放系统FreeFlow。考虑到完全去中心化社区治理通常还涉及治理代币、提案投票、内容治理和复杂的链上协同机制，而评论、检索索引、登录会话等高频交互数据若直接上链会带来较高的存储成本、确认延迟与维护复杂度，本文在系统取舍上采用Web2.5协同架构：使用Electron与React构建桌面端应用，使用Node.js、Express、Prisma和PostgreSQL构建链下业务服务，使用Solidity智能合约实现音乐资源发布、ERC-1155访问凭证、购买授权和收益分账，并通过Pinata/IPFS保存音频、封面和metadata资源。系统面向创作者提供作品草稿管理、素材上传、metadata生成、链上发布和访问状态查询等功能，面向听众提供链上资源检索、详情查看、购买访问权、播放、评论、下载和链上音乐库管理等功能，从而在“链上可信记录 + 链下高效服务”的边界内实现面向音乐社区的基础互动闭环。

    为提升链上音乐资源的可检索性，本文进一步设计了链上音乐资源检索与排序模块。该模块综合考虑标题、艺术家、专辑、流派、描述等文本字段，以及合约地址、Token ID、交易哈希和IPFS CID等链上身份字段，通过查询归一化、候选召回、特征评分、模糊相似度、新鲜度和热度计算完成排序，并返回可解释的排序原因。本文的可复现实验验证主要聚焦于智能合约自动化测试和检索排序离线评估两个部分。结果表明，合约自动化测试4项全部通过；在包含56条资源和62个查询的评估集上，完整排序算法的MRR达到0.9758，nDCG@10达到0.9646，在500条候选规模下平均响应时间为36.09 ms。实验结果说明，本文所实现的混合架构能够支持从创作者发布到用户购买播放与基础互动的核心业务流程。
}

\newcommand{\keywordsCN}{
    链上音乐；区块链；IPFS；Web2.5；访问控制；检索排序
}

\newcommand{\abstractEN}{
    With the development of digital music distribution and blockchain technology, traditional music platforms still rely heavily on centralized services for creator publishing, copyright credentials, revenue distribution, and access control. To address these issues, this thesis designs and implements FreeFlow, an on-chain music publishing and playback system for music-community scenarios. Because a fully decentralized community platform would further require governance tokens, proposal voting, content governance, and large-scale on-chain coordination, while frequently updated data such as comments, search indexes, and login sessions would incur high storage cost, confirmation latency, and maintenance overhead if written directly on-chain, this work adopts a Web2.5 collaborative architecture. The system uses Electron and React to build the desktop client, Node.js, Express, Prisma, and PostgreSQL to build the off-chain service layer, Solidity smart contracts to implement music publishing, ERC-1155 access credentials, purchase authorization, and revenue splitting, and Pinata/IPFS to store audio, cover images, and metadata resources. It provides creators with draft management, asset uploading, metadata generation, on-chain publishing, and access status checking, while providing listeners with on-chain resource retrieval, detail browsing, access purchase, playback, comments, downloading, and on-chain music library management.

    To improve the retrievability of on-chain music resources, this thesis also designs an on-chain music resource retrieval and ranking module. The module considers both textual fields such as title, artist, album, genre, and description, and on-chain identity fields such as contract address, token ID, transaction hash, and IPFS CID. It performs ranking through query normalization, candidate retrieval, feature scoring, fuzzy similarity, freshness, and popularity calculation, and returns explainable ranking reasons. The reproducible evaluation in this thesis focuses on smart-contract automation tests and offline retrieval experiments. The results show that all four contract tests pass, the full ranking algorithm reaches an MRR of 0.9758 and an nDCG@10 of 0.9646 on a benchmark with 56 resources and 62 queries, and the average latency remains 36.09 ms with 500 candidates. These results indicate that the hybrid architecture implemented in this work can support the core workflow from creator publishing to listener purchase, playback, and basic interaction.
}

\newcommand{\keywordsEN}{
    On-chain music; Blockchain; IPFS; Web2.5; Access control; Retrieval ranking
}
